\documentclass[10pt,a4paper]{amsart}
\usepackage{amsmath,amssymb,amsthm,graphicx}
% For commutative diagrams we use Xy-pic
\usepackage[all]{xy}

% We use 2cell for 2-commutative diagrams.
\xyoption{2cell}
% Theorem environments.
%
\theoremstyle{plain}
\newtheorem{theorem}[subsection]{Theorem}
\newtheorem{proposition}[subsection]{Proposition}
\newtheorem{lemma}[subsection]{Lemma}

\theoremstyle{definition}
\newtheorem{definition}[subsection]{Definition}
\newtheorem{example}[subsection]{Example}
\newtheorem{exercise}[subsection]{Exercise}
\newtheorem{situation}[subsection]{Situation}

\theoremstyle{remark}
\newtheorem{remark}[subsection]{Remark}
\newtheorem{remarks}[subsection]{Remarks}

\numberwithin{equation}{subsection}

% Macros
%
\def\lim{\mathop{\mathrm{lim}}\nolimits}
\def\colim{\mathop{\mathrm{colim}}\nolimits}
\def\Spec{\mathop{\mathrm{Spec}}}
\def\Hom{\mathop{\mathrm{Hom}}\nolimits}
\def\Ext{\mathop{\mathrm{Ext}}\nolimits}
\def\SheafHom{\mathop{\mathcal{H}\!\mathit{om}}\nolimits}
\def\SheafExt{\mathop{\mathcal{E}\!\mathit{xt}}\nolimits}
\def\Sch{\mathit{Sch}}
\def\Mor{\mathop{\mathrm{Mor}}\nolimits}
\def\Ob{\mathop{\mathrm{Ob}}\nolimits}
\def\Sh{\mathop{\mathit{Sh}}\nolimits}
\def\NL{\mathop{N\!L}\nolimits}
\def\CH{\mathop{\mathrm{CH}}\nolimits}
\def\proetale{{pro\text{-}\acute{e}tale}}
\def\etale{{\acute{e}tale}}
\def\QCoh{\mathit{QCoh}}
\def\Ker{\mathop{\mathrm{Ker}}}
\def\Im{\mathop{\mathrm{Im}}}
\def\Coker{\mathop{\mathrm{Coker}}}
\def\Coim{\mathop{\mathrm{Coim}}}

% Boxtimes
%
\DeclareMathSymbol{\boxtimes}{\mathbin}{AMSa}{"02}

%
% Macros for moduli stacks/spaces
%
\def\QCohstack{\mathcal{QC}\!\mathit{oh}}
\def\Cohstack{\mathcal{C}\!\mathit{oh}}
\def\Spacesstack{\mathcal{S}\!\mathit{paces}}
\def\Quotfunctor{\mathrm{Quot}}
\def\Hilbfunctor{\mathrm{Hilb}}
\def\Curvesstack{\mathcal{C}\!\mathit{urves}}
\def\Polarizedstack{\mathcal{P}\!\mathit{olarized}}
\def\Complexesstack{\mathcal{C}\!\mathit{omplexes}}
% \Pic is the operator that assigns to X its picard group, usage \Pic(X)
% \Picardstack_{X/B} denotes the Picard stack of X over B
% \Picardfunctor_{X/B} denotes the Picard functor of X over B
\def\Pic{\mathop{\mathrm{Pic}}\nolimits}
\def\Picardstack{\mathcal{P}\!\mathit{ic}}
\def\Picardfunctor{\mathrm{Pic}}
\def\Deformationcategory{\mathcal{D}\!\mathit{ef}}
\newenvironment{reference}{}{}
\newenvironment{slogan}{}{}
\newenvironment{history}{}{}
\usepackage{cleveref}
\crefname{originallemma}{original lemma}{original lemmas}
\Crefname{originallemma}{Original lemma}{Original lemmas}
\makeatletter
\newlabel{lemma-divisor-special-fiber}{{0C5Z}{source}{}{sourcepointer-1}{}}
\newlabel{lemma-divisor-special-fiber@cref}{{[originallemma][0][]0C5Z}{[1][0][]source}{}{}{}}
\makeatother
\crefname{originallemma}{original lemma}{original lemmas}
\Crefname{originallemma}{Original lemma}{Original lemmas}
\makeatletter
\newlabel{morphisms-lemma-characterize-proper}{{0BX5}{source}{}{sourcepointer-2}{}}
\newlabel{morphisms-lemma-characterize-proper@cref}{{[originallemma][0][]0BX5}{[1][0][]source}{}{}{}}
\makeatother
\crefname{originalsection}{original section}{original sections}
\Crefname{originalsection}{Original section}{Original sections}
\makeatletter
\newlabel{schemes-section-points}{{01J5}{source}{}{sourcepointer-3}{}}
\newlabel{schemes-section-points@cref}{{[originalsection][0][]01J5}{[1][0][]source}{}{}{}}
\makeatother
\crefname{originallemma}{original lemma}{original lemmas}
\Crefname{originallemma}{Original lemma}{Original lemmas}
\makeatletter
\newlabel{algebra-lemma-regular-ring-CM}{{00NQ}{source}{}{sourcepointer-4}{}}
\newlabel{algebra-lemma-regular-ring-CM@cref}{{[originallemma][0][]00NQ}{[1][0][]source}{}{}{}}
\makeatother
\crefname{originallemma}{original lemma}{original lemmas}
\Crefname{originallemma}{Original lemma}{Original lemmas}
\makeatletter
\newlabel{algebra-lemma-separable-smooth}{{00TV}{source}{}{sourcepointer-5}{}}
\newlabel{algebra-lemma-separable-smooth@cref}{{[originallemma][0][]00TV}{[1][0][]source}{}{}{}}
\makeatother
\crefname{originallemma}{original lemma}{original lemmas}
\Crefname{originallemma}{Original lemma}{Original lemmas}
\makeatletter
\newlabel{varieties-lemma-regular-functions-proper-variety}{{04L2}{source}{}{sourcepointer-6}{}}
\newlabel{varieties-lemma-regular-functions-proper-variety@cref}{{[originallemma][0][]04L2}{[1][0][]source}{}{}{}}
\makeatother
\crefname{originallemma}{original lemma}{original lemmas}
\Crefname{originallemma}{Original lemma}{Original lemmas}
\makeatletter
\newlabel{lemma-regular-model-field}{{0C68}{source}{}{sourcepointer-7}{}}
\newlabel{lemma-regular-model-field@cref}{{[originallemma][0][]0C68}{[1][0][]source}{}{}{}}
\makeatother
\crefname{originallemma}{original lemma}{original lemmas}
\Crefname{originallemma}{Original lemma}{Original lemmas}
\makeatletter
\newlabel{lemma-regular-model-genus}{{0C69}{source}{}{sourcepointer-8}{}}
\newlabel{lemma-regular-model-genus@cref}{{[originallemma][0][]0C69}{[1][0][]source}{}{}{}}
\makeatother
\title{A rational point on the curve produces a smooth rational point on the special fibre and forces its residue field, multiplicity and genus}
\author{The Stacks Project Authors}
\date{stacks-project:models.tex:local-snapshot-20261012; GNU FDL 1.2 or later}
\begin{document}
\maketitle
\noindent\textit{Editorial note.} The counted unit is the Stacks Project Lemma lemma-numerical-type-rational-point at source lines 5332-5344 together with its proof at source lines 5346-5372, a straight adjacent pairing in the chapter file. The chapter is the upstream file models.tex; the pool directory carries the longer name stacks-models and is not used as the chapter name anywhere in this card.\par\medskip
\noindent\textit{Editorial note.} One context span carries the regular-model situation at lines 4628-4637, which the retained statement names and which its own reference to the divisor lemma makes resolvable through a pointer; it is context rather than a counted unit and does not overlap the counted interval.\par\medskip
\noindent\textit{Editorial note.} The retained proof names two lemmas of this chapter and five results of the Algebra, Morphisms, Schemes and Varieties chapters. All are declared as non-load-bearing navigation pointers with their stable Stacks tags looked up in the shared tag index; the chapter ones are counted by cards of this lane and none of their intervals is reproduced here.\par\medskip
\noindent\textit{Editorial note.} The retained proof span starts at line 5346, the line of the opening marker read from the source, and contains exactly one proof environment.\par\medskip
\noindent\textit{Editorial note.} No citation occurs in the retained spans, so no bibliography entry is supplied.\par\medskip
\noindent\textit{Allowed standard prerequisites (editorial).}\begin{itemize}
\item The regular-model situation, smoothness at a point and the residue field of a component. The situation fixes a discrete valuation ring with residue field, a smooth projective curve over its fraction field, a regular proper model and the components of the special fibre with their multiplicities; a rational point of the curve extends to a morphism from the spectrum of the ring because the model is proper over it, and the image of the closed point lies on the special fibre. At that image the local ring of the special fibre is obtained by dividing out the uniformizer, its regularity follows from the regularity of the local ring of the model together with the Cohen-Macaulay property, and smoothness at the point follows from regularity combined with separability of the residue extension. The author's own lemmas of the chapter and the quoted results of the Algebra, Morphisms, Schemes and Varieties chapters are the allowed inputs.. The counted lemma is the computation that the chapter performs whenever the curve has a rational point: it pins down the numerical type of the component through that point and gives the equality of genera used later in the genus formula.
\end{itemize}\medskip
\ifdefined\hypertarget\hypertarget{sourcepointer-1}{}\fi
\noindent\textit{Source reference.} Original lemma 0C5Z: focused-seeds/AGP/sources/stacks-snapshot-20261011/tags.txt, source lines 15279--15279. Navigation pointer carried by the retained situation, which names this lemma when it writes the special fibre as a sum of its components with multiplicities; that lemma is counted by a separate card of this lane and its interval is not reproduced here. The number is the stable Stacks tag of the target, looked up in the shared tag index.\par\medskip
\ifdefined\hypertarget\hypertarget{sourcepointer-2}{}\fi
\noindent\textit{Source reference.} Original lemma 0BX5: focused-seeds/AGP/sources/stacks-snapshot-20261011/tags.txt, source lines 14968--14968. Allowed input of the Morphisms chapter quoted by the retained proof for the valuative criterion of properness, which extends the rational point; the pointer is registered so the reference resolves and the cited statement is not reproduced inside the span. The number is the stable Stacks tag of the target, looked up in the shared tag index.\par\medskip
\ifdefined\hypertarget\hypertarget{sourcepointer-3}{}\fi
\noindent\textit{Source reference.} Original section 01J5: focused-seeds/AGP/sources/stacks-snapshot-20261011/tags.txt, source lines 2002--2002. Allowed input of the Schemes chapter quoted by the retained proof for the description of a morphism out of the spectrum as a local ring homomorphism at the image point; the pointer is registered so the reference resolves and the cited material is not reproduced inside the span. The number is the stable Stacks tag of the target, looked up in the shared tag index.\par\medskip
\ifdefined\hypertarget\hypertarget{sourcepointer-4}{}\fi
\noindent\textit{Source reference.} Original lemma 00NQ: focused-seeds/AGP/sources/stacks-snapshot-20261011/tags.txt, source lines 868--868. Allowed input of the Algebra chapter quoted by the retained proof for the Cohen-Macaulay property of a regular local ring, which makes the quotient by the uniformizer regular; the pointer is registered so the reference resolves and the cited statement is not reproduced inside the span. The number is the stable Stacks tag of the target, looked up in the shared tag index.\par\medskip
\ifdefined\hypertarget\hypertarget{sourcepointer-5}{}\fi
\noindent\textit{Source reference.} Original lemma 00TV: focused-seeds/AGP/sources/stacks-snapshot-20261011/tags.txt, source lines 1089--1089. Allowed input of the Algebra chapter quoted by the retained proof for the passage from regularity with separable residue extension to smoothness; the pointer is registered so the reference resolves and the cited statement is not reproduced inside the span. The number is the stable Stacks tag of the target, looked up in the shared tag index.\par\medskip
\ifdefined\hypertarget\hypertarget{sourcepointer-6}{}\fi
\noindent\textit{Source reference.} Original lemma 04L2: focused-seeds/AGP/sources/stacks-snapshot-20261011/tags.txt, source lines 5993--5993. Allowed input of the Varieties chapter quoted by the retained proof for the identification of the residue field of a component with its global sections; the pointer is registered so the reference resolves and the cited statement is not reproduced inside the span. The number is the stable Stacks tag of the target, looked up in the shared tag index.\par\medskip
\ifdefined\hypertarget\hypertarget{sourcepointer-7}{}\fi
\noindent\textit{Source reference.} Original lemma 0C68: focused-seeds/AGP/sources/stacks-snapshot-20261011/tags.txt, source lines 15288--15288. Navigation pointer invoked by the retained proof for the result that the global sections of the special fibre form a field extension of the residue field; that lemma is counted by a separate card of this lane and its interval is not reproduced here. The number is the stable Stacks tag of the target, looked up in the shared tag index.\par\medskip
\ifdefined\hypertarget\hypertarget{sourcepointer-8}{}\fi
\noindent\textit{Source reference.} Original lemma 0C69: focused-seeds/AGP/sources/stacks-snapshot-20261011/tags.txt, source lines 15289--15289. Navigation pointer invoked by the retained proof for the equality between the genus of the curve and the genus read off from the numerical type of the special fibre; the cited statement is an allowed input of the same section and its proof is not reproduced here. The number is the stable Stacks tag of the target, looked up in the shared tag index.\par\medskip
\begin{situation}
\label{situation-regular-model}
Let $R$ be a discrete valuation ring with fraction field $K$,
residue field $k$, and uniformizer $\pi$.
Let $C$ be a smooth projective curve over $K$ with $H^0(C, \mathcal{O}_C) = K$.
Let $X$ be a regular proper model of $C$.
Let $C_1, \ldots, C_n$ be the irreducible components of the special
fibre $X_k$. Write $X_k = \sum m_i C_i$ as in
Lemma \ref{lemma-divisor-special-fiber}.
\end{situation}

\begin{lemma}
\label{lemma-numerical-type-rational-point}
In Situation \ref{situation-regular-model} assume $C$ has a $K$-rational point.
Then
\begin{enumerate}
\item $X_k$ has a $k$-rational point $x$ which is a smooth point of $X_k$
over $k$,
\item if $x \in C_i$, then $H^0(C_i, \mathcal{O}_{C_i}) = k$ and
$m_i = 1$, and
\item $H^0(X_k, \mathcal{O}_{X_k}) = k$ and $X_k$ has genus equal to
the genus of $C$.
\end{enumerate}
\end{lemma}

\begin{proof}
Since $X \to \Spec(R)$ is proper, the $K$-rational point extends to
a morphism $a : \Spec(R) \to X$ by the valuative criterion of properness
(Morphisms, Lemma \ref{morphisms-lemma-characterize-proper}).
Let $x \in X$ be the image under $a$ of the closed point of $\Spec(R)$.
Then $a$ corresponds to an $R$-algebra homomorphism
$\psi : \mathcal{O}_{X, x} \to R$
(see Schemes, Section \ref{schemes-section-points}).
It follows that $\pi \not \in \mathfrak m_x^2$ (since the image
of $\pi$ in $R$ is not in $\mathfrak m_R^2$).
Hence $\mathcal{O}_{X_k, x} = \mathcal{O}_{X, x}/\pi \mathcal{O}_{X, x}$
is regular (Algebra, Lemma \ref{algebra-lemma-regular-ring-CM}).
Then $X_k \to \Spec(k)$ is smooth at $x$ by
Algebra, Lemma \ref{algebra-lemma-separable-smooth}.
It follows that $x$ is contained in a unique irreducible component
$C_i$ of $X_k$, that $\mathcal{O}_{C_i, x} = \mathcal{O}_{X_k, x}$,
and that $m_i = 1$. The fact that $C_i$ has a
$k$-rational point implies that the field
$\kappa_i = H^0(C_i, \mathcal{O}_{C_i})$
(Varieties, Lemma \ref{varieties-lemma-regular-functions-proper-variety})
is equal to $k$. This proves (1). We have
$H^0(X_k, \mathcal{O}_{X_k}) = k$
because $H^0(X_k, \mathcal{O}_{X_k})$ is a field
extension of $k$ (Lemma \ref{lemma-regular-model-field})
which maps to $H^0(C_i, \mathcal{O}_{C_i}) = k$.
The genus equality follows from Lemma \ref{lemma-regular-model-genus}.
\end{proof}

\end{document}
